% !TeX root = part4.tex

\documentclass[../final.tex]{subfiles}

\begin{document}

% При отдельной сборке продолжаем нумерацию полного конспекта.
\ifSubfilesClassLoaded{\setcounter{chapter}{2}\setcounter{section}{2}}{}


% ============================================================
% Лекция 15.10
% ============================================================


% ============================================================
\section{Нерелятивистская теория возмущений}
% ============================================================

\begin{rbox}[left=18pt,right=18pt,top=14pt,bottom=10pt,unbreakable]{[]}
    \centering

    \begin{tikzpicture}[
        >=Latex,
        every node/.style={text=rtext}
    ]
        \draw[draw=rc4, line width=1.1pt, fill=rc4, fill opacity=0.08]
        (0,0) circle (1.25);

        \draw[->, draw=rc3, line width=1.3pt]
        (-4.2,-1.05) -- (-0.6,0.05);

        \draw[draw=rc3, line width=1.3pt]
        (-0.6,0.05) -- (0,0.25);

        \draw[->, draw=rc3, line width=1.3pt]
        (0,0.25) -- (3.9,-0.95);

        \fill[rc4] (0,0.25) circle (1.7pt);

        \node[font=\small] at (-3.5,-0.45) {$H_0$};
        \node[font=\small] at (0.05,-0.55) {$H_0+\hat V$};
        \node[font=\small] at (3.4,-0.25) {$H_0$};
    \end{tikzpicture}

    \captionsetup{
        font=footnotesize,
        labelfont=bf,
        skip=5pt
    }

    \captionof{figure}{
        Постановка задачи рассеяния. Вне области взаимодействия
        частица свободна и описывается гамильтонианом $H_0$,
        внутри действует возмущение $\hat V$.
    }
    \label{fig:perturbation_region}
\end{rbox}

% ------------------------------------------------------------
\subsection{Разложение по невозмущённым состояниям}
% ------------------------------------------------------------

Пусть собственные функции невозмущённого гамильтониана известны:

\begin{equation*}
    H_0\varphi_n=E_n\varphi_n,
    \qquad
    \int\varphi_i^*\varphi_j\,d^3x=\delta_{ij}.
\end{equation*}

В области взаимодействия

\begin{equation*}
    \left(H_0+\hat V(\mathbf x,t)\right)\psi
    =
    i\frac{\partial\psi}{\partial t}.
\end{equation*}

Решение ищем в виде разложения по полному набору $\varphi_n$
с зависящими от времени коэффициентами:

\begin{equation*}
    \psi(\mathbf x,t)
    =
    \sum_n a_n(t)\,\varphi_n(\mathbf x)\,e^{-iE_nt}.
\end{equation*}

Подставляя, получаем

\begin{equation*}
    i\sum_n
    \frac{da_n}{dt}\,
    \varphi_n(\mathbf x)\,e^{-iE_nt}
    =
    \sum_n
    \hat V(\mathbf x,t)\,a_n(t)\,\varphi_n(\mathbf x)\,e^{-iE_nt}.
\end{equation*}

Умножаем на $\varphi_f^*e^{iE_ft}$ и интегрируем по $d^3x$.
В силу ортонормированности слева остаётся одно слагаемое:

\begin{equation*}
    \frac{da_f}{dt}
    =
    -i\sum_n a_n(t)
    \int\varphi_f^*\hat V\varphi_n\,d^3x\;
    e^{i(E_f-E_n)t}.
\end{equation*}


% ------------------------------------------------------------
\subsection{Первый порядок теории возмущений}
% ------------------------------------------------------------

\rassumption{Предположение:} в момент $t_0=-T/2$ до взаимодействия
система находится в состоянии $i$:

\begin{equation*}
    a_i(t_0)=1,
    \qquad
    a_n(t_0)=0,\quad n\neq i.
\end{equation*}

Если возмущение мало, то в первом порядке в правой части можно
оставить только $a_i\simeq1$:

\begin{equation*}
    a_f(t)
    =
    -i\int_{-T/2}^{t}dt'
    \underbrace{
        \int\varphi_f^*\hat V\varphi_i\,d^3x
    }_{V_{fi}(t')}
    \,e^{i(E_f-E_i)t'}.
\end{equation*}

\rconcept{Амплитуда перехода} --- значение коэффициента после
окончания взаимодействия:

\begin{equation*}
    T_{fi}
    =
    a_f\!\left(t\to+\infty\right)
    =
    -i\int d^4x\;
    \varphi_f^*(x)\,\hat V(x)\,\varphi_i(x),
\end{equation*}

где теперь $\varphi_n(x)=\varphi_n(\mathbf x)e^{-iE_nt}$ включает
временную зависимость. Именно в такой ковариантной форме
амплитуда будет использоваться дальше.


% ------------------------------------------------------------
\subsection{Золотое правило Ферми}
% ------------------------------------------------------------

Для не зависящего от времени возмущения

\begin{equation*}
    T_{fi}
    =
    -iV_{fi}
    \int_{-T/2}^{T/2}
    e^{i(E_f-E_i)t}\,dt
    \;\xrightarrow[T\to\infty]{}\;
    -2\pi i\,V_{fi}\,\delta(E_f-E_i).
\end{equation*}

Вероятность перехода в единицу времени

\begin{equation*}
    W_{fi}
    =
    \lim_{T\to\infty}
    \frac{|T_{fi}|^2}{T}.
\end{equation*}

При возведении в квадрат один из интегралов даёт
$2\pi\delta(E_f-E_i)$, а второй при $E_f=E_i$ равен просто $T$:

\begin{equation*}
    |T_{fi}|^2
    =
    |V_{fi}|^2
    \left|
        \int_{-T/2}^{T/2}e^{i(E_f-E_i)t}\,dt
    \right|^2
    =
    |V_{fi}|^2\,
    T\,
    2\pi\delta(E_f-E_i).
\end{equation*}

Следовательно,

\begin{equation*}
    W_{fi}
    =
    2\pi\,|V_{fi}|^2\,\delta(E_f-E_i).
\end{equation*}

Переход в одно точное состояние физического смысла не имеет:
нужно проинтегрировать по группе конечных состояний с
\rconcept{плотностью конечных состояний} $\rho(E_f)$.
Получаем \rconcept{золотое правило Ферми}:

\begin{equation*}
    W_{fi}
    =
    2\pi\,|V_{fi}|^2\,\rho(E_f).
\end{equation*}


% ============================================================
\section{Электромагнитное взаимодействие}
% ============================================================

% ------------------------------------------------------------
\subsection{Минимальная подстановка}
% ------------------------------------------------------------

Как включить взаимодействие с электромагнитным полем?
Как и в классической механике, для частицы с зарядом $-e$
(электрон) делается замена

\begin{equation*}
    \mathbf p\;\longrightarrow\;\mathbf p+e\mathbf A,
    \qquad
    \hat p^\mu\;\longrightarrow\;\hat p^\mu+eA^\mu,
    \qquad
    A^\mu=(\varphi,\mathbf A),
\end{equation*}

где $A^\mu$ --- \rnote{четырёхмерный потенциал}, а не поле.
Поскольку $\hat p^\mu=i\partial^\mu$, это то же самое, что

\begin{equation*}
    \partial^\mu
    \;\longrightarrow\;
    \partial^\mu-ieA^\mu.
\end{equation*}

Уравнение Клейна--Гордона

\begin{equation*}
    \left(\partial_\mu\partial^\mu+m^2\right)\varphi=0
\end{equation*}

переходит в \rconcept{модифицированное уравнение}

\begin{equation*}
    \left(
        \partial^\mu\partial_\mu
        -ie\,\partial_\mu A^\mu
        -ie\,A^\mu\partial_\mu
        -e^2A^2
        +m^2
    \right)\varphi
    =
    0,
\end{equation*}

которое удобно записать как

\begin{equation*}
    \left(\partial^\mu\partial_\mu+m^2+\hat V\right)\varphi=0,
    \qquad
    \hat V
    =
    -ie\left(\partial_\mu A^\mu+A^\mu\partial_\mu\right)
    -e^2A^2.
\end{equation*}

Здесь $\partial_\mu A^\mu$ --- оператор, действующий на всё,
что стоит справа, т.\,е. $\partial_\mu(A^\mu\varphi)$.


% ------------------------------------------------------------
\subsection{Ток перехода}
% ------------------------------------------------------------

Подставим $\hat V$ в амплитуду перехода. Слагаемое
$e^2A_\mu A^\mu$ \rassumption{отбрасываем}, так как это
следующий порядок по $e$:

\begin{equation*}
    T_{fi}
    =
    -i\int\varphi_f^*(x)\,\hat V(x)\,\varphi_i(x)\,d^4x
    =
    -e\int\varphi_f^*
    \left(
        A^\mu\partial_\mu+\partial_\mu A^\mu
    \right)
    \varphi_i\,d^4x.
\end{equation*}

Во втором слагаемом интегрируем по частям
(поля на бесконечности обращаются в нуль):

\begin{equation*}
    \int\varphi_f^*\,\partial_\mu\!\left(A^\mu\varphi_i\right)d^4x
    =
    -\int\left(\partial_\mu\varphi_f^*\right)A^\mu\varphi_i\,d^4x.
\end{equation*}

Тогда

\begin{equation*}
    T_{fi}
    =
    -e\int A^\mu
    \left[
        \varphi_f^*\left(\partial_\mu\varphi_i\right)
        -
        \varphi_i\left(\partial_\mu\varphi_f^*\right)
    \right]
    d^4x
    =
    -i\int j^{fi}_\mu\,A^\mu\,d^4x,
\end{equation*}

где

\begin{equation*}
    j^{fi}_\mu
    =
    -ie
    \left(
        \varphi_f^*\,\partial_\mu\varphi_i
        -
        \varphi_i\,\partial_\mu\varphi_f^*
    \right)
\end{equation*}

--- \rconcept{ток перехода} $i\to f$. Он совпадает с током
$-e\,j^\mu$ (заряд электрона $-e$) из уравнения Клейна--Гордона,
если заменить $\varphi^*\to\varphi_f^*$, $\varphi\to\varphi_i$.

Для плоских волн $\varphi=Ne^{-ipx}$:

\begin{equation*}
    j^{fi}_\mu
    =
    -eN_iN_f\left(p_i+p_f\right)_\mu e^{i(p_f-p_i)x},
\end{equation*}

\begin{equation*}
    T_{fi}
    =
    ieN_iN_f\left(p_i+p_f\right)_\mu
    \int A^\mu\,e^{iqx}\,d^4x,
    \qquad
    q=p_f-p_i.
\end{equation*}

Величина $q$ --- переданный 4-импульс.


% ============================================================
\section{Рассеяние на кулоновском центре}
% ============================================================

% ------------------------------------------------------------
\subsection{Система Хевисайда--Лоренца}
% ------------------------------------------------------------

В КЭД принято использовать
\rassumption{рационализированную систему Хевисайда--Лоренца},
в которой

\begin{equation*}
    F=\frac{q_1q_2}{4\pi r^2},
    \qquad
    \operatorname{div}\mathbf E=\rho,
    \qquad
    \varphi=\frac{q}{4\pi r}.
\end{equation*}

Постоянная тонкой структуры в ней

\begin{equation*}
    \alpha=\frac{e^2}{4\pi}\approx\frac{1}{137},
    \qquad
    r_e=\frac{e^2}{4\pi m_e}=\frac{\alpha}{m_e}.
\end{equation*}


% ------------------------------------------------------------
\subsection{Амплитуда рассеяния}
% ------------------------------------------------------------

\begin{rbox}[left=18pt,right=18pt,top=14pt,bottom=10pt,unbreakable]{[]}
    \centering

    \begin{tikzpicture}[
        >=Latex,
        every node/.style={text=rtext},
        photon/.style={
            draw=rtext,
            line width=0.9pt,
            decorate,
            decoration={snake, amplitude=1.6pt, segment length=6pt}
        }
    ]
        \coordinate (V) at (0,0);

        \draw[->, draw=rc3, line width=1.2pt] (-2.6,1.7) -- (-1.1,0.72);
        \draw[draw=rc3, line width=1.2pt] (-1.1,0.72) -- (V);
        \draw[->, draw=rc3, line width=1.2pt] (V) -- (2.6,1.7);

        \node[font=\small] at (-2.2,0.95) {$p_i$};
        \node[font=\small] at (2.2,0.95) {$p_f$};

        \draw[photon] (V) -- (0,-1.9);
        \node[font=\small] at (0.45,-0.95) {$q$};

        \draw[draw=rc4, line width=1.1pt] (-0.2,-2.1) -- (0.2,-1.7);
        \draw[draw=rc4, line width=1.1pt] (-0.2,-1.7) -- (0.2,-2.1);
        \node[font=\small, right=6pt] at (0.2,-1.9) {заряд $e$};

        \fill[rc4] (V) circle (1.7pt);
    \end{tikzpicture}

    \captionsetup{
        font=footnotesize,
        labelfont=bf,
        skip=5pt
    }

    \captionof{figure}{
        Рассеяние заряженной частицы на статическом кулоновском
        центре. Частица передаёт импульс $q=p_f-p_i$
        через виртуальный фотон.
    }
    \label{fig:coulomb_center}
\end{rbox}

Возьмём статический кулоновский потенциал

\begin{equation*}
    A^\mu
    =
    \left(
        \frac{e}{4\pi|\mathbf r|},\,
        \mathbf 0
    \right).
\end{equation*}

В $T_{fi}$ остаётся только $\mu=0$ и $(p_i+p_f)_0=E_i+E_f$:

\begin{equation*}
    T_{fi}
    =
    ieN_iN_f\left(E_i+E_f\right)
    \int e^{i(E_f-E_i)t}\,dt
    \int\frac{e}{4\pi|\mathbf r|}\,
    e^{-i\mathbf q\mathbf r}\,d^3x.
\end{equation*}

Интеграл по времени даёт $2\pi\delta(E_f-E_i)$ ---
рассеяние упругое. Пространственный интеграл --- фурье-образ
кулоновского потенциала. Направим ось $Oz$ вдоль $\mathbf q$:

\begin{equation*}
    \frac{e}{4\pi}
    \int\frac{1}{r}\,
    e^{-iqr\cos\theta}\,
    r^2\,dr\,d\varphi\,d\cos\theta
    =
    \frac{e}{|\mathbf q|^2}.
\end{equation*}

\remark[Как считать этот интеграл.]
Интеграл по $r$ сходится только условно. Его регуляризуют,
заменяя потенциал на экранированный (юкавский)
$\dfrac{1}{4\pi r}e^{-\mu r}$:

\begin{equation*}
    \int\frac{e^{-\mu r}}{4\pi r}\,e^{-i\mathbf q\mathbf r}\,d^3x
    =
    \frac{1}{q}\int_0^\infty e^{-\mu r}\sin(qr)\,dr
    =
    \frac{1}{|\mathbf q|^2+\mu^2},
\end{equation*}

а затем полагают $\mu\to0$. На лекции было отмечено, что
потенциал вида
$A^0\to\dfrac{1}{4\pi r}e^{-mr}$ \rnote{тоже имеет место быть}
--- это поле, переносимое массивной частицей; он даёт
$1/(|\mathbf q|^2+m^2)$ вместо $1/|\mathbf q|^2$.
Пока такие поля не рассматриваем.

Для упругого рассеяния 4-вектор $q=(0,\mathbf q)$, поэтому
$q^2=-|\mathbf q|^2$. Итого

\begin{equation*}
    T_{fi}
    =
    -iN_iN_f\,\mathcal M\,2\pi\delta(E_f-E_i),
    \qquad
    \mathcal M
    =
    -e^2\frac{E_i+E_f}{|\mathbf q|^2}.
\end{equation*}

Здесь $\mathcal M$ --- \rconcept{матричный элемент},
не содержащий нормировок и $\delta$-функции.


% ------------------------------------------------------------
\subsection{Сечение: формула Резерфорда}
% ------------------------------------------------------------

Вероятность перехода в единицу времени

\begin{equation*}
    W_{fi}
    =
    e^4N_i^2N_f^2\,
    \frac{\left(E_i+E_f\right)^2}{|\mathbf q|^4}\,
    2\pi\delta(E_f-E_i)
    =
    2\pi N_i^2N_f^2\,|\mathcal M|^2\,\delta(E_f-E_i).
\end{equation*}

Сечение --- это вероятность, делённая на поток падающих частиц
и умноженная на число конечных состояний:

\begin{equation*}
    d\sigma
    =
    \frac{W_{fi}}{J}\,\rho_f.
\end{equation*}

\rassumption{Нормировка:} одна частица в объёме $V$, т.\,е.
$\rho=2E|N|^2=1/V$, $N^2=1/(2EV)$. Тогда поток $J=v_i/V$,
а число конечных состояний

\begin{equation*}
    \rho_f
    =
    \frac{V\,d^3p_f}{(2\pi\hbar)^3}.
\end{equation*}

Объём $V$ сокращается, и

\begin{equation*}
    d\sigma
    =
    \frac{e^4}{v_i}\,
    \frac{1}{2E_i}\,
    \frac{1}{2E_f}\,
    \frac{\left(E_i+E_f\right)^2}{|\mathbf q|^4}\,
    2\pi\delta(E_f-E_i)\,
    \frac{d^3p_f}{(2\pi)^3}.
\end{equation*}

Поскольку $p\,dp=E\,dE$, то $d^3p_f=p_f^2\,dp_f\,d\Omega=p_fE_f\,dE_f\,d\Omega$:

\begin{equation*}
    d\sigma
    =
    \frac{e^4}{v_i}\,
    \frac{1}{2E_i}\,
    \frac{1}{2E_f}\,
    \frac{\left(E_i+E_f\right)^2}{|\mathbf q|^4}\,
    2\pi\delta(E_f-E_i)\,
    \frac{p_fE_f\,dE_f\,d\Omega}{(2\pi)^3}.
\end{equation*}

После интегрирования по $dE_f$ положим $E_i=E_f=E$,
$|\mathbf p_i|=|\mathbf p_f|=p$, $v=p/E$. Переданный импульс
выражается через угол рассеяния:

\begin{equation*}
    |\mathbf q|^2
    =
    |\mathbf p_f-\mathbf p_i|^2
    =
    4p^2\sin^2\frac{\theta}{2}.
\end{equation*}

В результате

\begin{equation*}
    \frac{d\sigma}{d\Omega}
    =
    \frac{e^4E^2}{4\pi^2|\mathbf q|^4}
    =
    \frac{\alpha^2}{4p^2v^2\sin^4\dfrac{\theta}{2}},
\end{equation*}

или, вводя $r_e=\dfrac{e^2}{4\pi m_e}$,

\begin{equation*}
    \frac{d\sigma}{d\Omega}
    =
    \frac{r_e^2}{4}
    \left(
        \frac{m_e}{\beta p}
    \right)^2
    \frac{1}{\sin^4\dfrac{\theta}{2}}.
\end{equation*}

Это \rresult{формула Резерфорда} --- та самая, что использовалась
в первой главе при выводе ионизационных потерь.
Квантовый расчёт в первом порядке совпадает с классическим.


% ------------------------------------------------------------
\subsection{Рассеяние двух частиц: первый взгляд}
% ------------------------------------------------------------

Если вместо статического центра взять вторую движущуюся заряженную
частицу, например в процессе

\begin{equation*}
    \pi^++K^+\rightarrow\pi^++K^+,
\end{equation*}

то каждая частица создаёт свой ток, а фотон с импульсом $q$
переносит взаимодействие между ними:

\begin{rbox}[left=18pt,right=18pt,top=14pt,bottom=10pt,unbreakable]{[]}
    \centering

    \begin{tikzpicture}[
        >=Latex,
        every node/.style={text=rtext},
        photon/.style={
            draw=rtext,
            line width=0.9pt,
            decorate,
            decoration={snake, amplitude=1.6pt, segment length=6pt}
        },
        scalar/.style={draw=rc3, line width=1.2pt}
    ]
        \coordinate (A) at (0,1);
        \coordinate (B) at (0,-1);

        \draw[scalar, ->] (-2.3,2.4) -- (-1.15,1.7);
        \draw[scalar] (-1.15,1.7) -- (A);
        \draw[scalar, ->] (A) -- (2.3,2.4);

        \draw[scalar, ->] (-2.3,-2.4) -- (-1.15,-1.7);
        \draw[scalar] (-1.15,-1.7) -- (B);
        \draw[scalar, ->] (B) -- (2.3,-2.4);

        \draw[photon] (A) -- (B);

        \node[font=\small] at (-2.5,2.1) {$p_1$};
        \node[font=\small] at (2.5,2.1) {$p_2$};
        \node[font=\small] at (-2.5,-2.1) {$p_3$};
        \node[font=\small] at (2.5,-2.1) {$p_4$};
        \node[font=\small] at (-2.95,2.65) {$\pi^+$};
        \node[font=\small] at (-2.95,-2.65) {$K^+$};
        \node[font=\small] at (0.4,0) {$q$};

        \fill[rc4] (A) circle (1.7pt);
        \fill[rc4] (B) circle (1.7pt);
    \end{tikzpicture}

    \captionsetup{
        font=footnotesize,
        labelfont=bf,
        skip=5pt
    }

    \captionof{figure}{
        Рассеяние $\pi^+K^+\rightarrow\pi^+K^+$ за счёт обмена
        виртуальным фотоном.
    }
    \label{fig:piK_first}
\end{rbox}

\begin{equation*}
    \mathcal M
    \sim
    e\left(p_1+p_2\right)^\mu
    \left(-\frac{1}{q^2}\right)
    e\left(p_3+p_4\right)_\mu.
\end{equation*}

Множитель $-1/q^2$ --- это фурье-образ кулоновского поля
(\rconcept{пропагатор фотона}). Подробно этот процесс разобран
в следующем разделе.

\end{document}
